\newpage

\newvideofile{Mastrom}{Maschenstromverfahren}
\begin{frame}
    
	\ftx{Learning objectives Mesh flow method}
	\begin{Key point}{mesh flow method}
        The Students 
        \begin{itemize}
			\item are familiar with the steps of the mesh flow process.
            \item can analyse electrical networks using the mesh current method. 
        \end{itemize}                   
    \end{Key point}

	\speech{MastromLernz}{1}{Neben dem Knotenpotentialverfahren bietet das Maschenstromverfahren eine weitere Möglichkeit, 
    um elektrische Netzwerke zu analysieren. 
    Die für die Durchführung des Maschenstromverfahrens erforderlichen Schritte werden folgend vorgestellt.} 
\end{frame}

\begin{frame}
	\renewcommand{\figurename}{Figure}
    \fta{Mesh flow method}

    \s{
        In addition to node potential analysis for investigating electrical networks, the mesh current method can also be used for calculation. 
        In the mesh current method, a system of equations 
        is set up with which the mesh currents can be determined. In order to set up the system of equations for the mesh current method, 
        some preparations must be made. In the following steps, the electrical 
        network from Figure \ref{BildMaschenstromverfahren} is to be calculated using the mesh current method.
 
    }

    \b{
		\begin{itemize}
			\item Preparation of the network
			\item Define meshes
			\item Determine resistance matrix
			\item Assign source voltages
			\item set up a system of equationsn
		\end{itemize}
	}

	\fu{   
        \resizebox{0.7\textwidth}{!}{%                                         
        \input{Tikz/src/Netzwerk_Maschenstrom}
        }                                                  
    }{{\bf Network for the mesh flow method.} Electrical network with a voltage source, a current source and three resistors. The network is used to 
    explain the mesh current method.. \label{BildMaschenstromverfahren}}   
    
    \speech{MastromSchritte}{1}{Zu beginn sollte wie auch beim Knotenpotentialverfahren das elektrische Netzwerk vorbereitet werden.
    Das beinhaltet unter anderem das Zusammenfassen von Reihen und Parallelschaltungen wie auch das eventuelle Umwandeln der Quellen. 
    Folgend werden die Maschen und die zugehörigen Maschenströme festgelegt, die Widerstandsmatrix erstellt, die Quellspannungen zugeordnet 
    und daraus ein Gleichungssystem aufgestellt.} 
\end{frame}




\begin{frame}
    \renewcommand{\figurename}{Figure}
\ftb{Preparation of the network}

    \s{
        The mesh current method works with voltage sources and resistance values. For this purpose, any current sources must be transformed into voltage sources.
        This type of conversion of a current source into a voltage source is shown once again in Figure \ref{BildMasch1}.
        Here, the current source and the parallel internal resistance are transformed into a voltage source with a series internal resistance. 
        Here, the current source and the parallel 
        internal resistance are converted into a voltage source with a series internal resistance. 
    }

	\b{
		\begin{itemize}
			\item Conversion of current sources into voltage sources:
		\end{itemize}
	}    

    \fu{                                     
        \input{Tikz/src/Umwandlung_Stromquelle}                                                   
    }{{\bf Conversion of a power source.} Conversion of a current source with parallel resistance to a voltage source with series resistance. \label{BildMasch1}}   

    \s{
        Equation \ref{GleichungUmrechnungStromquelle} shows the calculation of the voltage value of the voltage source 
        from the specified current value of the current source and the resistance value of the internal resistance. 
    }

    \onslide<2->{
	\b{
		\begin{itemize}
			\item Conversion of current values to voltage values:
		\end{itemize}
	}

    \begin{eq}
         U_\mathrm{q} = R_\mathrm{i} \cdot I_\mathrm{q} \label{GleichungUmrechnungStromquelle}
    \end{eq}}

    \s{
        In addition, possible conductance values of electrical resistors must be converted into resistance values. Equation 
        \ref{GleichungUmrechnungLeitwert} shows how the resistance value is determined from the reciprocal of the conductance. 
    }

    \onslide<3->{
	\b{
		\begin{itemize}
			\item Calculation of conductance values:
		\end{itemize}
	}

    \begin{eq}               
         R_\mathrm{i} = \frac{1}{G_\mathrm{i}} \label{GleichungUmrechnungLeitwert}
    \end{eq}}

    \speech{MastromVorb}{1}{Eventuelle Stromquellen müssen für das Maschenstromverfahren in Spannungsquellen umgewandelt werden.
    Wie hier dargestellt werden die Stromquellen mit Parallelwiderstand in Spannungsquellen mit Reihenwiderstand umgewandelt.} 
    \speech{MastromVorb}{2}{Die Stromwerte der Stromquellen werden über das ohmsche Gesetz in Spannungswerte überführt.} 
    \speech{MastromVorb}{3}{Falls Leitwerte anstatt von Widerstandswerten angegeben werden, können diese über den Kehrwert ermittelt werden,
    da beim Maschenstromverfahren mit den Widerstandswerten gearbeitet wird.} 
\end{frame}

\begin{frame}
    \renewcommand{\figurename}{Figure}
    \ftx{Preparation of the network}

    \s{
        In Figure \ref{BildMasch2}, the current source of the network shown is converted into a voltage source.
        The parallel connection of the current source and resistor $R_2$ becomes a series connection consisting
        of the new voltage source and resistor $R_2$.
    }

    \b{
        Conversion of the current source into a voltage source:
        \begin{itemize}
            \item Replace the power source with a voltage source
            \item The parallel resistor becomes a series resistor.
        \end{itemize}
    }

    \fu{                                     % Beginn Bild Netzwerk 1 mit rotem Kreis
		\resizebox{\textwidth}{!}{%
        \input{Tikz/src/Umwandlung_Stromquelle_Netzwerk}
        }                                           
    }{{\bf Conversion of a power source in an electrical network.} Conversion of the current source with parallel resistance 
    of the network to be calculated into a voltage source with series resistance. The transformation of the source takes place in the red 
    circle. \label{BildMasch2}}  
    
    \speech{MastromUmw}{1}{Die Umwandlung der Stromquelle in eine Spannungsquelle wird für das angegebene Netzwerk dargestellt.
    Aus dem Parallelwiderstand der Stromquelle wird der Reihenwiderstand für die Spannungsquelle.} 
\end{frame}

\begin{frame}
    \ftx{Mesh flow method}

    \begin{Key point}{Mesh flow method}
        For the mesh current method, voltage sources and resistance values are required to analyse an electrical network.
    \end{Key point}

    \speech{MastromMrk}{1}{Für das Maschenstromverfahren werden zur Analyse eines elektrischen Netzwerkes 
    Spannungsquellen und Widerstandswerte benötigt.} 
\end{frame}


\begin{frame}
    \renewcommand{\figurename}{Figure}
\ftb{Define meshes}

    \s{
        For mesh current analysis, the meshes must be defined. The mesh currents are determined via the meshes.
        In addition, the resistances are assigned to the meshes. Figure \ref{BildMasch3} shows 
        the network presented with the meshes $M_1$ and $M_2$ marked.
    } 
    
    \b{
		\begin{itemize}
			\item Determination of meshes and mesh flows
		\end{itemize}
	}

    \fu{    
        \resizebox{0.7\textwidth}{!}{%                                                 
        \input{Tikz/src/Maschen_im_Netzwerk}
        }                                                            
    }{{\bf Mesh in the electrical network.} Determination of meshes $M_1$ and $M_2$ and the associated mesh flows.  
    \label{BildMasch3}}                                                

    \s{
        
        The mesh current $I_\mathrm{M1}$ flows through all components located at mesh $M_1$. Similarly, 
        all components at mesh $M_2$ are traversed by the mesh current $I_\mathrm{M2}$. The two mesh currents
        $I_\mathrm{M1}$ and $I_\mathrm{M2}$ are noted vectorially in equation \ref{GleichungMaschenströme}. 
    }

    \begin{eq}
        I_\mathrm{M} = 
        \begin{bmatrix}
            I_\mathrm{M1} \\
            \\
            I_\mathrm{M2} \\
        \end{bmatrix} \label{GleichungMaschenströme}
    \end{eq}      

    \speech{MastromMasch}{1}{Die Maschen M 1 und M 2 werden entsprechend definiert. 
    Ihnen zugehörig werden die Maschenströme IM eins und IM zwei der beiden Maschen im Vektor der Maschenströme IM festgehalten.}     
\end{frame}


\begin{frame}
    \renewcommand{\figurename}{Figure}
    \ftb{Determine resistance matrix}

    \s{
        In mesh $M_1$, the mesh current $I_{M1}$ encounters the two resistors $R_0$ and $R_1$. Mesh $M_1$ is shown in 
        blue in Figure \ref{BildWiderstandsmatrix}. The mesh current $I_{M2}$ of mesh $M_2$ flows through the 
        resistors $R_1$ and $R_2$. Resistor $R_1$ is therefore part of both meshes $M_1$ and $M_2$ and is traversed by both 
        mesh currents. If a resistor is traversed by several mesh currents, it is considered a coupling resistor 
        between the meshes.
        In the resistance matrix according to equation
        \ref{GleichungWiderstandsmatrix}, the resistances of the meshes and the coupling resistances are assigned. 
        Here the loop resistances from the meshes are entered in the main diagonal. This means that the series connection $R_0 + R_1$ 
        becomes the element of the first column and row for the first mesh. The series connection $R_1 + R_2$ is entered in the element 
        of the second column and row for the second mesh. The coupling resistances between the respective meshes are entered in the remaining elements of the resistance matrix. 
        The coupling resistance is given a positive 
        sign if both mesh currents flow in the same direction and a negative sign if they flow in opposite directions.
        In alternating current technology, the resistance matrix is also referred to as the mesh impedance matrix. 
    }

    \b{
		\begin{itemize}
			\item<2-> Main diagonal with the resistance values directly adjacent to the meshes
			\item<3-> Other elements with resistance values that lie directly between the meshes involved
		\end{itemize}
	}

    \fu{    
        \resizebox{0.7\textwidth}{!}{%                                                 
        \input{Tikz/src/Festlegung_der_Zweige}
        }                                                            
    }{{\bf Determination of the branches.} Electrical network with highlighted branches for creating the resistance matrix for the mesh current method
    \label{BildWiderstandsmatrix}}

    \b{
		\only<1>{
		\begin{eq}
			R_\mathrm{M} =
			\begin{bmatrix} 
				0 & 0 \\
				\\
				0 & 0 \\
			\end{bmatrix} 
		\end{eq}
		}
		\only<2>{
		\begin{eq}
			R_\mathrm{M} =
			\begin{bmatrix} 
				\color{blue}{R_0 + R_1} & 0 \\
				\\
				0 & R_1 + R_2 \\
			\end{bmatrix} 
		\end{eq}
		}
		\only<3>{
		\begin{eq}
			R_\mathrm{M} =
			\begin{bmatrix} 
				\color{blue}{R_0 + R_1} & \color{green}{-R_1} \\
				\\
				\color{green}{-R_1} & R_1 + R_2 \\
			\end{bmatrix} 
		\end{eq}
		}
	}

    \s{
    \begin{eq}
        R_\mathrm{M} =
        \begin{bmatrix} 
            \color{blue}{R_0 + R_1} & \color{green}{-R_1} \\
			\\
			\color{green}{-R_1} & R_1 + R_2 \\
        \end{bmatrix} \label{GleichungWiderstandsmatrix}
    \end{eq}
    }

    \speech{MastromWiderstands}{1}{Als nächstes wird die Widerstandsmatrix für das angegebene elektrische Netzwerk bestimmt.}
    \speech{MastromWiderstands}{2}{Die an den Maschen beteiligten Widerstände werden auf der Hauptdiagonalen der Matrix eingetragen.
    Für die Masche M 1 sind das die in blau gekennzeichneten Widerstände R null und R eins.
    An der Masche M 2 sind die Widerstände R eins und R 2 beteiligt. } 
    \speech{MastromWiderstands}{3}{Die Widerstände welche an zwei oder mehr Maschen beteiligt sind, 
    werden zwischen den Maschen eingetragen. 
    In diesem Fall handelt es sich hier um den Widerstand R eins. 
    Die Widerstände zwischen den Maschen werden mit einem negativen Vorzeichen versehen.}  
\end{frame}


\begin{frame}
    \ftb{Assign source voltages}

    \s{
        The mesh must be assigned the stress sources, known as source stresses. The mesh $M_1$ has the 
        source stress $U_0$. The mesh $M_2$ has the source stress $U_1$. If the direction of the voltage arrow 
        and the mesh direction are the same, the source voltages have a negative sign. If the directions of the voltage arrow and
        the mesh circulation are opposite, the sign of the source voltage is positive. If there is no source in a mesh, a $0$ is entered.
    }

    \b{
        \fu{    
        \resizebox{0.7\textwidth}{!}{%                                                 
        \input{Tikz/src/Maschen_im_Netzwerk}
        }                                                            
    }    
    }
   
    \begin{eq}
        U_\mathrm{M} = 
        \begin{bmatrix}
            U_0 \\
            \\
            -U_1 \\
        \end{bmatrix} \label{GleichungQuellspannungen}
    \end{eq}

    \speech{MastromQuell}{1}{In dem Vektor U M werden die Quellströme festgehalten. 
    An der Masche M 1 befindet sich die Spannungsquelle U null. 
    An der Masche M 2 befindet sich die Spannungsquelle U eins.
    Verläuft der Spannungspfeil der Spanungsquelle mit dem Maschenumlauf,
    wird ein negatives Vorzeichen, 
    bei einem entgegengesetzten Spannungspfeil ein positives Vorzeichen gewählt.}  
\end{frame}


\begin{frame}
\ftb{Set up a system of equations}

\s{
    According to Ohm's law, electrical voltage is calculated as the product of electrical resistance and current.
    In mesh current analysis, the resistance matrix $R_\mathrm{M}$ and the vector of source voltages $U_\mathrm{M}$ are used to calculate the vector of mesh currents $I_\mathrm{M}$.
    The system of equations is represented in Equation \ref{GleichungMasche1}.  
    }

\begin{eq}
\begin{bmatrix} 
    R_0 + R_1 & -R_1 \\
    \\
    -R_1 & R_1 + R_2 \\
\end{bmatrix}   
\cdot
\begin{bmatrix}
    I_\mathrm{M1} \\
    \\
    I_\mathrm{M2} \\
\end{bmatrix}
=
\begin{bmatrix}
    U_0 \\
    \\
    -U_1 \\
\end{bmatrix} \label{GleichungMasche1}
\end{eq}

\s{
    The quantities sought are calculated from the mesh currents. For example, the currents $I_0$, $I_1$ and $I_2$ are sought.
    The results of the mesh currents are equivalent to the currents that flow through a single mesh current. 
    In the electrical network shown, the current of mesh $M_1$ is $I_\mathrm{M1} = I_0$ and the current of mesh $M_2$ is $I_\mathrm{M2} = -I_2$. The
    current through resistor $R_1$ can be calculated using the node equation: $I_1 = I_0 + I_2$.
}

\speech{MastromGleich}{1}{Nun kann aus der Widerstandsmatrix, 
dem Vektor der Maschenströme 
und dem Vektor der Quellspannungen 
das notwendige Gleichungssystem aufgestellt und gelöst werden.
Aus den so zu bestimmenden Maschenströmen können anschließend die gesuchten Ströme bestimmt werden.
Hier ist beispielsweise der Strom I 0 gleich dem Maschenstrom I emm eins.
Der Maschenstrom I emm 2 ist gleich dem negativen Wert von I zwei, 
da die Verläufe der beiden Ströme entgegengesetzt sind.
Zum Schluss kann der Strom I eins über die Knotenregel aus den Strömen I 0 und I 2 berechnet werden.
Somit sind dann alle Größen des Netzwerkes identifiziert worden.}
\end{frame}
\newpage


%%-------------------------NETZWERK2------------------------------%%



\begin{frame}
    \ftx{Example: Mesh flow method}
    \begin{expl}{Mesh flow method}{}

    \s{
        Analysis of the electrical network from Figure \ref{BildMasch4} using mesh current analysis. 
    }

    \begin{enumerate}
        \only<1>{	
            \item[]

        The following steps should be carried out for mesh flow analysis:

    \begin{itemize}
        \item[a)] Preparation of the network
		\item[b)] Define meshes and mesh flows
		\item[c)] Determine resistance matrix
		\item[d)] Assign source voltages
		\item[e)] Set up a system of equations
    \end{itemize}

    \fu{  
        \resizebox{0.4\textwidth}{!}{%                                     
        \input{Tikz/src/Beispiel_Maschenanalyse}
        }                  
    }{{\bf Example.} Mesh current analysis on an electrical network. \label{BildMasch4}}                                            
    
        }


        \only<2>{			
            \item[a)] Preparation of the network
			\begin{figure}[H]
                \resizebox{0.7\textwidth}{!}{%
                    \input{Tikz/src/Vorbereitung_Netzwerk}	
                    }\centering                                                           
                \end{figure}
            }
        \only<3>{
            \item[b)] Define meshes and mesh flows
            \onslide<3>{
                \begin{figure}[H]
                \resizebox{0.4\textwidth}{!}{%                                                        
                    \input{Tikz/src/Maschen_definieren}
                    }\centering                                                           
                \end{figure}
            \pause
            
            \begin{eq}
                I_\mathrm{M} = 
                \begin{bmatrix}  
                    I_\mathrm{M1} \\
                    \\
                    I_\mathrm{M2} \\
                    \\
                    I_\mathrm{M3} \\ 
                \end{bmatrix}\nonumber
            \end{eq}

            }
        }
        \only<4>{
            \item[c)] Determine resistance matrix
            \onslide<4>{
                \begin{eq}
                    R_\mathrm{M} =
                    \begin{bmatrix}                    
                        R_0 + R_3 + R_4 & -R_3 & -R_4 \\
                        \\
                        -R_3 & R_1 + R_3 + R_5 & -R_5\\
                        \\
                        -R_4 & -R_5 & R_2 + R_4 + R_5 \\
                    \end{bmatrix} \nonumber
                \end{eq} 
            }
        }
        \only<5>{
            \item[d)] Assign source voltages
            \onslide<5>{
                \begin{eq}
                    U_\mathrm{M} = 
                    \begin{bmatrix} 
                        U_0 \\
                        \\
                        U_1 \\
                        \\
                        -U_2\\
                    \end{bmatrix} \nonumber
                \end{eq} 
            }
        }
        \only<6>{
            \item[e)] Set up a system of equations
            \onslide<6>{
                \begin{eq}
                    \begin{bmatrix} 
                        R_0 + R_3 + R_4 & -R_3 & -R_4 \\
                        \\
                        -R_3 & R_1 + R_3 + R_5 & -R_5\\
                        \\
                        -R_4 & -R_5 & R_2 + R_4 + R_5 \\
                    \end{bmatrix}   
                    \cdot
                    \begin{bmatrix}
                        I_\mathrm{M1} \\
                        \\
                        I_\mathrm{M2} \\
                        \\
                        I_\mathrm{M3} \\
                    \end{bmatrix}
                    =
                    \begin{bmatrix}
                        U_0 \\
                        \\
                        U_1 \\
                        \\
                        -U_2\\
                    \end{bmatrix}\nonumber
                \end{eq} 
                \pause
                Calculate currents: 
                \begin{eqa}   
                    I_3 = I_0 - I_1\nonumber\\
                    I_4 = I_3 - I_5\nonumber\\
                    I_0 = I_4 - I_2\nonumber\\
                    -I_5 = I_2 + I_1\nonumber
                \end{eqa}  
            }
        }
    \end{enumerate}
\end{expl}{}{}

\speech{MastromBsp}{1}{Die Schritte zur Anwendung des Maschenstromverfahrens wurden nun erläutert. 
Folgend wird anhand des abgebildeten Beispielnetzwerkes noch einmal die grundsätzliche Arbeitsweise wiederholt.
Hierzu wird das Netzwerk enstprechend für das Maschenstromverfahren vorbereitet, die Maschen und Maschenströme definiert,
die Widerstandsmatrix bestimmt, die Quellspannungen zugeordnet und anschließend das Gleichungssystem aufgestellt.}
\speech{MastromBsp}{2}{Bei dem Beispielnetzwerk müssen keine Widerstände zusammengefasst werden.
Jedoch wird die Stromquelle mit Parallelwiderstand in eine Spannungsquelle mit Serienwiderstand umgewandelt.}
\speech{MastromBsp}{3}{Für das Beispielnetzwerk werden drei Maschen M 1, M 2 und M 3 mit ihren zugehörigen 
Maschenströmen I emm eins, I emm zwei und I emm drei definiert.}
\speech{MastromBsp}{4}{An der Masche M 1 sind die Widerstände R null, R drei und R 4 beteiligt. Die Widerstände R eins, R drei und R fünf liegen an der Masche M 2. Die Masche M 3 verfügt über die Widerstände R zwei, R vier und R fünf. 
Zwischen den Maschen M eins und M zwei liegt der Widerstand R drei. Zwischen den Maschen M eins und M 3 liegt der Widerstand R vier und zwischen den Maschen M 2 und M 3 liegt der Widerstand R fünf. Die Widerstände werden in die Widerstandsmatrix R emm eingetragen.}
\speech{MastromBsp}{5}{Der Vektor der Quellspannungen ergibt sich nach U emm.
Für die Masche M 1 wird die Spannungsquelle U null aufgeschrieben.
Die Spannungsquelle U eins ist an der Masche M 2 beteiligt 
und für die Masche M 3 wird die Spannungsquelle minus U zwei notiert.}
\speech{MastromBsp}{6}{Mit der Widerstandsmatrix, dem Vektor der Maschenströme 
und dem Vektor der Quellspannungen kann das Gleichungssystem aufgestellt werden.
Nach dem Lösen der Gleichung sind die Maschenströme bekannt und können den Strömen zugeordnet werden.
Hier entspricht der Maschenstrom I emm 1 dem Strom I null.
Der Maschenstrom I emm 2 entspricht dem Strom I eins
und der Maschenstrom I emm 3 entspricht dem Strom minus I zwei.
Die übrigen Ströme des Netzwerkes können anschließend über die Knotenregel bestimmt werden,
bis letztendlich alle Größen des Netzwerkes bekannt sein sollten.}
\end{frame}
